\subsection{Subgroup Vulnerability and System Resilience}
\label{subsec:subgroup_resilience}

The impact of behavioral delay aversion ($\kappa$) is not distributed uniformly across the population, leading to significant variations in epidemic outcomes and system resilience. Figure~\ref{fig:vulnerability_dashboard} quantifies these disparities by comparing representative epidemic trajectories and final cumulative attack rates across eight socio-demographic and behavioral subgroups.

As illustrated in Panel A of Figure~\ref{fig:vulnerability_dashboard}, behavioral delay functions as a "force multiplier" for transmission. Compared to the \textit{Supply-only (Ideal)} baseline (final attack rate of 22.46\%), the inclusion of behavioral friction for the \textit{Overall Sample} ($\kappa = 0.225$) shifts the epidemic curve significantly upward, leading to a projected final attack rate of 63.58\%. For the most resistant subgroup---\textit{Neutral Institutional Trust} ($\kappa = 0.356$)---the system faces a near-total collapse of preventive efficacy, with the infection peak arriving earlier and the final attack rate reaching 65.88\%.

Panel B reveals a clear hierarchy of vulnerability. Subgroups traditionally considered marginalized or high-risk exhibit the highest structural hurdles to timely vaccination. Specifically, individuals with \textit{Primary Education} (65.75\% attack rate), \textit{Adults Aged 65+} (65.38\%), and those with a \textit{Prior Vaccine Reaction} (65.17\%) all face infection risks substantially higher than the population average. Notably, the elevated $\kappa$ among \textit{Adults Aged 65+} suggests that biological vulnerability is compounded by behavioral resistance to administrative delay, creating a "double burden" of risk.

These findings suggest that "system resilience" against an infectious outbreak is fundamentally constrained by behavioral parameters. A vaccination strategy that assumes uniform uptake timing will systematically fail the most vulnerable segments of society. To improve resilience, policy interventions must move beyond aggregate supply metrics and focus on reducing the specific behavioral "time-tax" faced by high-$\kappa$ subgroups, particularly through targeted trust-building and streamlined delivery for the elderly and those with chronic medical conditions.
